\chapter{LITERATURE REVIEW}

This chapter presents a review of references related to the Inventory Management System. The reviewed works cover trading business workflow, spreadsheet-based business records, web-based inventory systems, inventory management concepts, and AI-assisted query support. These topics are important because the project is not only a stock-counting system. It also connects supplier purchasing, customer sales, quotations, receivables, payables, dashboard summaries, and a limited read-only AI assistant.

\section{North American Industry Classification System: Wholesale Trade \cite{naics_2022_wholesale_trade}}
The North American Industry Classification System describes wholesale trade as an intermediate step in the distribution of merchandise. Wholesalers are organized to sell or arrange the purchase and sale of goods for resale, durable non-consumer goods, and raw or intermediate materials. This reference supports the business background of this project because the system is designed for a trading business that works between suppliers and customers.

In this project, the same idea appears in the main workflow. The business buys products from suppliers, stores received stock, and sells products to customers. Therefore, the system needs to manage both sides of the transaction. Supplier data, purchase records, customer data, sales records, and stock availability must be connected so that the user can understand the full flow of goods and money.

\section{Qtier-Rapor: Managing Spreadsheet Systems and Improving Corporate Performance, Compliance and Governance \cite{bishop2008spreadsheet}}
This paper discusses how businesses often use spreadsheet systems to fill gaps in their core administrative systems. Spreadsheets are flexible and easy to start with, but the paper explains that larger spreadsheet systems can become harder to manage and control. This is relevant to small and medium-sized businesses because many businesses begin with simple spreadsheets before moving to a more structured system.

The problem is similar to the motivation of this project. A trading business may be able to record products, purchases, sales, and payments in spreadsheets at first. However, when the number of products, suppliers, customers, and transaction statuses increases, the records become harder to check. This supports the need for a web application with a database, backend validation, and clearer workflow separation.

\section{Sole Traders Still Not Prepared for Digital Tax Returns in April \cite{times_sage_sole_traders_2025}}
This article reports research from Sage about sole traders and their preparation for digital tax reporting. It shows that many small business owners still use older methods such as spreadsheets, bank statements, and pen-and-paper records for financial administration. Although the article focuses on tax reporting, it is useful for this project because it shows that manual and semi-manual record keeping is still common in small business environments.

For an inventory trading business, the same issue can affect daily operations. If stock, purchasing, sales, billing, and payment information are kept in separate files or folders, it becomes difficult to answer simple questions quickly. Examples include checking which products are low in stock, which sales can be billed, and which purchases are ready for payment. This project addresses that issue by keeping related records in one web-based system.

\section{A Full-Stack Web Solution for Online FMCG Management System Using MERN: Design, Implementation, and Evaluation \cite{paper_1}}
This study presents the development of a full-stack web solution for managing fast-moving consumer goods. It discusses a web-based system with inventory management, order processing, stock updates, billing, and user roles. The reference is related to this project because both systems focus on using a web application to reduce manual business work and improve access to operational records.

The main difference is the technology stack and project scope. The reviewed system uses the MERN stack, while this project uses React for the frontend, Django REST Framework for the backend, and MySQL for the relational database. This project also focuses specifically on SME trading workflows, including purchases, sales, quotations, billing notes, payment batches, credit notes, stock validation, dashboard summaries, and limited AI-assisted operation.

\section{Essentials of Inventory Management \cite{paper_2}}
This book explains that inventory management is more than counting available stock. It includes demand, replenishment, physical stock handling, inventory value, and the business decisions connected to stock. These ideas support the design of the inventory and dashboard features in this project.

In the implemented system, stock is calculated from received purchase quantities and committed sale quantities. The dashboard and inventory page also show reorder-related information, pending demand, pending purchases, average cost, stock value, and product movement. These features follow the general inventory management idea that a business should understand both the physical quantity of stock and the financial meaning of that stock.

\section{Aisha: A Custom AI Library Chatbot Using the ChatGPT API \cite{paper_3}}
This article describes the development of a custom chatbot using the ChatGPT API for a university library. The chatbot was designed to answer user questions using prepared system information and a controlled application context. The reference is useful for this project because it shows how a chatbot can be connected to a specific domain instead of being used as a general-purpose conversation tool.

In this Inventory Management System, the AI assistant follows a limited and read-only scope. It can support questions about stock and fulfillment, customer or supplier summaries, receivables and payables, and reference or line-item lookup. It does not create, approve, edit, cancel, or delete records. This limitation is important because inventory and finance records must remain controlled by the main application workflow and backend validation.

\section{Summary}
The reviewed references support the main direction of this project. The wholesale trade reference explains why a trading business needs to manage the connection between suppliers, products, and customers. The spreadsheet and small-business references show why manual records, spreadsheets, and separated files can become difficult to manage as business operations grow.

The web-based inventory system reference supports the idea of building an integrated application for stock and transaction workflows. The inventory management book supports the calculation and reporting features, especially stock availability, replenishment, and inventory value. The chatbot reference supports the idea of adding a controlled AI-assisted interface, while also showing why the assistant should be limited to a clear domain.

Based on these works, this project focuses on developing a practical web-based inventory management system for SME trading businesses. The system combines React, Django REST Framework, and MySQL to manage business records in a structured way. It also adds dashboard summaries and a read-only AI assistant to help users find supported operational information more quickly.
